% ============================================================================
%  Dirichlet Character Atlas: A Weil-Kernel Cartography
%  of Low-Conductor Sector Gaps
%
%  Author: Lukas Geiger, Independent Researcher, Bernau
%  Draft v2 — revised after N=600 server analysis
%  Status: Final Draft
% ============================================================================
\documentclass[a4paper,10pt,twocolumn]{article}

\usepackage[a4paper,top=2.0cm,bottom=2.5cm,left=1.8cm,right=1.8cm,columnsep=0.6cm]{geometry}
\usepackage[utf8]{inputenc}
\usepackage[T1]{fontenc}
\usepackage{lmodern}
\usepackage{microtype}
\usepackage[english]{babel}
\usepackage{amsmath,amssymb,amsfonts,mathtools,bm}
\usepackage{graphicx,float,booktabs,array,multirow}
\usepackage{xcolor}
\usepackage{xurl}
\usepackage[unicode=true,colorlinks=true,linkcolor=blue!60!black,citecolor=blue!60!black,urlcolor=blue!60!black]{hyperref}
\usepackage[numbers,square,sort&compress]{natbib}
\bibliographystyle{unsrtnat}

\usepackage{titlesec}
\renewcommand{\thesection}{\Roman{section}}
\renewcommand{\thesubsection}{\Alph{subsection}}

\usepackage{amsthm}
\newtheorem{theorem}{Theorem}[section]
\newtheorem{proposition}[theorem]{Proposition}
\newtheorem{corollary}[theorem]{Corollary}
\newtheorem{lemma}[theorem]{Lemma}
\theoremstyle{definition}
\newtheorem{definition}[theorem]{Definition}
\newtheorem{example}[theorem]{Example}
\theoremstyle{remark}
\newtheorem{remark}[theorem]{Remark}

% --- Custom macros ---
\newcommand{\gap}{\operatorname{gap}}
\newcommand{\chifn}{\chi}
\newcommand{\Zmodq}{(\mathbb{Z}/q\mathbb{Z})^*}
\newcommand{\halfline}{\tfrac{1}{2}}
\newcommand{\Deltachi}{\Delta_{\chifn}}
\newcommand{\archdiff}{\operatorname{arch\_diff}}
\newcommand{\primediff}{\operatorname{prime\_diff}_{\mathrm{matrix}}}
\DeclareMathOperator{\sgn}{sgn}
\providecommand*{\HyperDestNameFilter}[1]{en.#1}

% ============================================================================
\title{The Dirichlet Character Atlas\\
  \large --- A Weil-Kernel Numerical Atlas for Dirichlet $L$-Function Sector Gaps\\
  \normalsize Galerkin Diagnostics and the Boundaries of Leading-Order Theory}
\author{Lukas Geiger\\\small Independent Researcher, Bernau}
\date{\today}

\hypersetup{
  pdftitle={The Dirichlet Character Atlas},
  pdfauthor={Lukas Geiger},
  pdfsubject={Weil-kernel numerical atlas for Dirichlet L-function sector gaps},
  pdfkeywords={Dirichlet L-functions, Weil quadratic form, sector gaps, Galerkin diagnostics, Hilbert-Polya}
}

\begin{document}
\twocolumn[
\maketitle
\begin{abstract}
\noindent
This paper is an instantiation of Functional Stability Theory
(FST) in the mathematical domain: it extends the three-component
decomposition $\mathrm{RH}(X) = \mathrm{GBC}(X) + C^{(2)}(X) +
\mathrm{Hurwitz}(X)$ (see Math-Master~\cite{MathMaster}, \S2.4) to
twisted $L$-functions, where the sector gap is the quantitative
signature of $C^{(2)}_{\chifn}$. The Math-Master
\cite{MathMaster} develops the overall zeta-zoo architecture and
the trinity UCU + SGE + Weil-form transversality; the Physics-Master
\cite{RFEPFramework} develops the companion physical formulation
via the Renormalized Free Energy Principle.

We develop a Weil-kernel numerical atlas of sector gaps for Dirichlet
$L$-functions associated to ten low-conductor real characters
$\chi_D$, $D\in\{5,8,12,13,17,21,24,29,33,60\}$, at Galerkin truncations
$N\in\{200,400,600\}$ and cut-off $\lambda=20000$. Within the
three-component decomposition the sector gap is the quantitative
signature of $C^{(2)}_{\chifn}$.
Our main results are \emph{negative but structurally informative}. First,
the leading-order approximation $\phi^+_{\chifn}\approx\phi^-_{\chifn}$
of the two sector ground states is empirically falsified: a systematic
16-window test, including the empirical ground state itself, fails the
success criterion $9/10$ sign accuracy, $R^2\geq 0.8$. Second, the
ostensibly ``exact'' formula
\[
  \gap_{\chifn} = -2\!\!\sum_{p,m}\!\!\frac{\chifn(p)^m\log p}{p^{m/2}}
    \Deltachi(m\log p) + \archdiff_{\chifn} + O(N^{-1})
\]
with
$\Deltachi(t):=(\phi^+_{\chifn}\ast\phi^+_{\chifn})(t)+(\phi^-_{\chifn}\ast\phi^-_{\chifn})(t)$
reduces, once both ground states are taken from the same Galerkin
diagonalization, to a matrix-decomposition tautology: an apparent
$9/9$ sign accuracy at $N=200$ disappears as a numerical interpolation
artifact when the analysis is repeated at $N=600$. Third, the
$\chi_{21}$ gap undergoes a sign flip between $N=200$ and $N=600$
($+0.282\to-0.004$), illustrating that even the Galerkin operator
itself is not informative below its convergence scale. We catalog the
ten characters along two filter constants (parity and root number, both
$+1$ in the sample) and two informative axes (archimedean factor weight
and gap signature), exhibit the systematic Siegel-Walfisz compression at
$N=600$, and position the result as a precise diagnostic of why
numerical Galerkin routes cannot, by themselves, resolve
$C^{(2)}_{\chifn}$. A genuine predictive formulation would require
Paley-Wiener asymptotic control of the sector ground states that
bypasses the Galerkin truncation; we close with a roadmap toward
such a branch-local Hilbert-P\'olya framework.
\end{abstract}
\vspace{1em}
]
\clearpage

\section{Introduction}
\label{sec:intro}

\subsection{Motivation}
\label{sec:intro-motivation}

The direct-frontier dominance approach to the Riemann hypothesis
\cite{RH_II} represents the even-parity sector of the Weil quadratic form
on $L^2[-L,L]$, $L=\log\lambda$, as dominating the odd-parity sector
uniformly in $\lambda$, at the rate
$\mathrm{gap}(\lambda)\gtrsim\sqrt{\lambda}$. The construction is
unconditional, purely variational, and uses only the three structural
ingredients of $\zeta(s)$: the Euler product, the Gamma factor, and
a single functional equation.

A natural and central question is whether the same architecture
transfers to twisted $L$-functions $L(s,\chi)$ for non-trivial
Dirichlet characters $\chi$, thereby addressing the Generalized
Riemann Hypothesis (GRH) one character at a time. In the meta-programme
of \cite{MetaGalerkinBridge}, the RH-type question for a zeta family
$X$ is diagnosed through a three-component decomposition
\begin{equation}
\label{eq:three-component}
  \mathrm{RH}_X
    \;=\;
  \mathrm{GBC}_X \;+\; C^{(2)}_X \;+\; \mathrm{Hurwitz}_X,
\end{equation}
where $\mathrm{GBC}_X$ captures the Galerkin bridge to a self-adjoint
operator, $\mathrm{Hurwitz}_X$ the reality-of-zeros transfer, and
$C^{(2)}_X$ the Connes-style operator identification at the heart of the
Hilbert-P\'olya programme. For the trivial character,
$C^{(2)}_{X=\zeta}$ is closed unconditionally by the even-dominance
theorem of \cite{RH_II}; for non-trivial $\chi$, $C^{(2)}_{\chifn}$ is
the hard residual component. The present paper is an atlas of what
Galerkin-numerical techniques can and, more importantly, cannot say
about this residual.

\subsection{What this paper does}
\label{sec:intro-contribution}

We catalogue the Galerkin sector gap
$\gap_{\chifn}(\lambda)
  :=\lambda_{\min}(Q^-_{\chifn})-\lambda_{\min}(Q^+_{\chifn})$
at $\lambda=20000$ for ten Kronecker-real characters
$\chi_D$, $D\in\{5,8,12,13,17,21,24,29,33,60\}$,
across three Galerkin resolutions $N\in\{200,400,600\}$. For each
character we compute both sector ground states separately
$\phi^\pm_{\chifn}$, the primitive Weil kernel prediction
\begin{equation}
\label{eq:predictor}
  S_{\chifn}(\lambda)
    = -2\sum_{\substack{p\nmid q\\m\log p\leq 2L}}
      \frac{\chifn(p)^m\log p}{p^{m/2}}\;\Deltachi(m\log p),
\end{equation}
with
$\Deltachi(t)
  =(\phi^+_{\chifn}\!*\phi^+_{\chifn})(t)
    +(\phi^-_{\chifn}\!*\phi^-_{\chifn})(t)$,
the archimedean contribution
$\archdiff_{\chifn}
  :=\langle\phi^-,K^{\mathrm{arch}}_-\phi^-\rangle
    -\langle\phi^+,K^{\mathrm{arch}}_+\phi^+\rangle$,
and the full Galerkin gap $\gap_{\mathrm{gal}}^{(N)}$. We position the
four quantities relative to each other, across all three Galerkin
resolutions, and relative to the empirical (highest-$N$) gaps as a
reference.

\subsection{What this paper finds}
\label{sec:intro-findings}

Our contribution is deliberately diagnostic. The key findings are:

\begin{enumerate}
\item[(F1)] \textbf{Falsification of the leading-order reduction.}
  Under the approximation
  $\phi^+_{\chifn}\approx\phi^-_{\chifn}\approx\phi_\infty$ used in
  \cite{RH_II} for the trivial character, the prime-side predictor
  reduces to a universal autocorrelation integral. A systematic test
  of 16 window models (including the empirical single-sector ground
  state) achieves at best $6/10$ sign accuracy and $R^2=0.15$.
  The leading-order reduction is false for non-trivial $\chi$.
\item[(F2)] \textbf{Matrix-decomposition tautology.}
  If \emph{both} sector ground states are taken from the same
  Galerkin diagonalization, the identity
  $\gap_{\mathrm{gal}}^{(N)} = S_{\chifn}+\archdiff_{\chifn}$
  is not a prediction but a consequence of the linearity of
  expectation on $W=W_{\mathrm{arch}}+W_{\mathrm{prime}}$. An apparent
  $9/9$ sign accuracy at $N=200$ (modulo the known $N$-oscillator
  $\chi_{21}$) dissolves under $N=600$ server-level validation: the
  predictor $S_{\chifn}+\archdiff_{\chifn}$ drops to $8/10$ owing
  to two sign failures: $\chi_{21}$ ($N$-oscillator,
  see (F3) and \S\ref{sec:chi21}) and $\chi_{29}$ (numerical
  interpolation artifact at a near-zero gap).
\item[(F3)] \textbf{$\chi_{21}$ as an $N$-oscillator.}
  At $N=200$ one has $\gap_{\mathrm{gal}}=+0.282$; at $N=600$ one has
  $\gap_{\mathrm{gal}}=-0.004$. The sign flip is not a formula defect
  but a Galerkin truncation effect for composite conductors that only
  settles above $N\sim 500$. This both demonstrates that the Galerkin
  operator is itself non-trivially $N$-dependent, and explains why a
  $9/9$ at $N=200$ was not surprising once $\chi_{21}$ was excluded:
  the excluded character was the very one whose $N$-dependence was not
  yet integrated into the test.
\item[(F4)] \textbf{Residual content: the Weil-kernel architecture is
  rigorous.} The sector-kernel decomposition
  (Theorem~\ref{thm:kernel}), the anti-diagonal sector difference
  (Theorem~\ref{thm:sector-diff}), and the spectral-domain form of the
  gap via the $L'/L$ function remain valid as structural statements
  independent of the leading-order reduction. We exhibit them in
  Section~\ref{sec:setup} and use them to classify the ten characters
  along four orthogonal axes in Section~\ref{sec:atlas}.
\end{enumerate}

Taken together, these findings deliver a precise diagnosis: the
$C^{(2)}_{\chifn}$-component of \eqref{eq:three-component}, for
non-trivial $\chi$, is not accessible to Galerkin-numerical closure
because the relevant signal lives in the difference
$\phi^+_{\chifn}-\phi^-_{\chifn}$, which is fine structure lost in
every natural leading-order approximation, and because the full
Galerkin operator is a tautology relative to its own eigenvectors.
A genuinely predictive formulation must identify
$\Deltachi$ analytically, e.g.\ via Paley-Wiener asymptotics
\cite{PaleyWiener} adapted from \cite{RH_II} to the
character-twisted setting.

\subsection{Organization}

Section~\ref{sec:setup} summarizes the Weil quadratic form and the
sector-kernel decomposition. Section~\ref{sec:formula} states the
gap formula and contrasts the reduced and unreduced versions.
Section~\ref{sec:validation} presents the numerical validation at
$N\in\{200,400,600\}$. Section~\ref{sec:atlas} catalogs the ten
characters along four axes. Section~\ref{sec:tautology} formulates
and proves the tautology. Section~\ref{sec:chi21} treats the
$\chi_{21}$ $N$-oscillation in detail. Section~\ref{sec:discussion}
discusses implications for \eqref{eq:three-component}. Section~\ref{sec:outlook}
outlines the Paley-Wiener roadmap and the branch-local
Hilbert-P\'olya program as the indicated next step.

\section{Setup: Weil quadratic form and sector kernels}
\label{sec:setup}

We summarize the structural framework used throughout; the full
derivations are given in \cite{AnalyticKernel}.

\subsection{\texorpdfstring{Dirichlet $L$-functions}{Dirichlet L-functions} and the Weil formula}

Let $\chifn$ be a primitive Dirichlet character mod $q$ with
$\chifn(-1)=+1$ (even). Set $a_{\chifn}=0$ and define the completed
$L$-function
\[
  \Lambda(s,\chifn)
    =\Bigl(\frac{q}{\pi}\Bigr)^{s/2}\Gamma\!\Bigl(\frac{s}{2}\Bigr)
      L(s,\chifn),
\]
satisfying the functional equation
$\Lambda(s,\chifn)=\epsilon_{\chifn}\Lambda(1-s,\bar\chifn)$ with
$|\epsilon_{\chifn}|=1$. For a real (i.e., Kronecker) character
$\bar\chifn=\chifn$ and $\epsilon_{\chifn}=+1$; we work exclusively
with the ten real characters $\chifn_D$ with $D\in\{5,8,12,13,17,21,
24,29,33,60\}$ in this atlas. The non-trivial zeros
$\rho=\halfline+i\gamma$ are assumed on the critical line throughout
(GRH for the tested conductors). The database pages of \cite{LMFDB}
provide computed low-lying zeros on the critical line for the
corresponding Dirichlet $L$-functions. We use these entries only as
finite numerical consistency data; they are not a certificate that all
zeros below a specified height have been verified on the critical line.
A theoretical proof of GRH for $L(s,\chifn_D)$ for any specific $D$ in
the sample is, of course, open.

For an even compactly supported test function $h\in C_c^\infty
(\mathbb{R})$, the Weil explicit formula
\cite{Weil1952,IwaniecKowalski} reads
\begin{equation}
\label{eq:weil}
  \sum_{\gamma}h(\gamma)
    = W^{\chifn}_\infty[h]
    - 2\sum_{\substack{p\nmid q\\m\geq 1}}
      \frac{\chifn(p)^m\log p}{p^{m/2}}\,\hat h(m\log p),
\end{equation}
where the sum on the left runs over all imaginary parts of non-trivial
zeros of $L(s,\chifn)$, and the archimedean weight is
\[
  W^{\chifn}_\infty[h]
    = \frac{1}{2\pi}\!\int_{-\infty}^{\infty}\! h(t)
      \Bigl[\log\!\frac{q}{\pi} + \Re\,\psi\!\Bigl(\tfrac14+\tfrac{it}{2}\Bigr)\Bigr]\,dt.
\]

\subsection{\texorpdfstring{Quadratic form on $H_L$}{Quadratic form on HL}}

Fix $\lambda>1$ and set $L:=\log\lambda$, $H_L:=L^2[-L,L]$. For a
real test function $f\in C_c^\infty([-L,L])$, put
$h(t):=2\pi|\hat f(t)|^2$. Then $h$ is non-negative even, and
\eqref{eq:weil} induces the \emph{Weil quadratic form}
\begin{equation}
\label{eq:weil-qf}
\begin{aligned}
  Q_{\chifn}[f]
    &:=\sum_{\gamma}2\pi|\hat f(\gamma)|^2 \\
    &=W^{\chifn}_\infty[2\pi|\hat f|^2] \\
    &\quad -2\sum_{p,m}\frac{\chifn(p)^m\log p}{p^{m/2}}
      \int f(y)f(y-m\log p)\,dy.
\end{aligned}
\end{equation}
$Q_{\chifn}$ admits a translation-invariant symbol: defining the
primitive convolution distribution
\begin{equation}
\label{eq:kappa}
\begin{aligned}
  \kappa_{\chifn}(u)
    &:=-\sum_{p\nmid q,\,m\geq 1}
      \frac{\chifn(p)^m\log p}{p^{m/2}} \\
    &\quad\times
      \bigl[\delta(u-m\log p)+\delta(u+m\log p)\bigr],
\end{aligned}
\end{equation}
and the archimedean kernel
$K^{\mathrm{arch}}_{\chifn}(u):=\mathcal{F}^{-1}[\rho_{\chifn}](u)$
with $\rho_{\chifn}(t)=\log(q/\pi)+\Re\,\psi(1/4+it/2)$, one obtains
the kernel
$K_{\chifn}(x,y)=\kappa_{\chifn}(x-y)+K^{\mathrm{arch}}_{\chifn}(x-y)$
satisfying $Q_{\chifn}[f]=\iint K_{\chifn}(x,y)f(x)f(y)\,dx\,dy$ on
$C_c^\infty([-L,L])$.

\subsection{Sector decomposition}

Let $\mathcal{P}:H_L\to H_L$, $(\mathcal{P}f)(x)=f(-x)$, denote the
reflection. $\mathcal{P}$ is unitary self-adjoint with $\mathcal{P}^2
=I$; the spectral projections $P^\pm=(I\pm\mathcal{P})/2$ split
$H_L=H_L^+\oplus H_L^-$ into even/odd sectors. Standard Fourier bases
are
\[
\begin{aligned}
  e_0^+&=\tfrac{1}{\sqrt{2L}}, &
  e_n^+(x)&=\tfrac{1}{\sqrt L}\cos(n\pi x/L),\\
  e_n^-(x)&=\tfrac{1}{\sqrt L}\sin(n\pi x/L).
\end{aligned}
\]
The Galerkin truncation at basis size $N$ uses modes $0\leq n\leq N$
in each sector.

\subsection{Sector kernels and the anti-diagonal difference}

Because $K_{\chifn}(x,y)=K_{\chifn}(-x,-y)$, the sector-projected
kernels take the form
\begin{theorem}[Sector kernels; {\cite{AnalyticKernel}, Thm.~3.2}]
\label{thm:kernel}
On $H_L^\pm$ respectively,
\[
\begin{aligned}
  K_{\chifn}^{\pm}(x,y)
    &=\tfrac{1}{2}\bigl[\kappa_{\chifn}(x-y)
       \pm\kappa_{\chifn}(x+y)\bigr]\\
    &\quad+\tfrac{1}{2}\bigl[K^{\mathrm{arch}}_{\chifn}(x-y)
       \pm K^{\mathrm{arch}}_{\chifn}(x+y)\bigr].
\end{aligned}
\]
\end{theorem}

\begin{proof}[Proof sketch]
The reflection $\mathcal{P}f(x):=f(-x)$ commutes with translation
$T_a$ via $\mathcal{P}T_a=T_{-a}\mathcal{P}$, and the kernel
$K_\chi(x,y)=\kappa_\chi(x-y)+K^{\mathrm{arch}}_\chi(x-y)$ is
translation-invariant. Symmetrising under $\mathcal{P}$:
$P^\pm K_\chi P^\pm=(K_\chi\pm\mathcal{P}K_\chi\mathcal{P})/2$,
where $\mathcal{P}K_\chi\mathcal{P}(x,y)=K_\chi(-x,-y)=K_\chi(x,y)$
by the parity of $\kappa_\chi$ and $K^{\mathrm{arch}}_\chi$ (both
even on $\mathbb{R}$). The off-diagonal symmetrisation
$P^\pm K_\chi(1-P^\pm)$ vanishes on the appropriate sector, leaving
the displayed expression. The action on the test-function space
$C_c^\infty([-L,L])$ is well-defined because $\kappa_\chi$ is a
tempered distribution supported on $\{u=\pm m\log p\}$, all of which
lie in $[-2L,2L]$; pairing with $f\in C_c^\infty([-L,L])$ produces a
finite sum over $(p,m)$ with $m\log p\leq 2L$. Full computational
details, including the verification that the symmetrisation
preserves self-adjointness, are in \cite{AnalyticKernel}, \S3.2.
\end{proof}

Projected eigenvalue problems
$K_{\chifn}^{\pm}\phi_{\chifn}^{\pm}=\lambda_1(\chifn,\pm)\phi_{\chifn}^{\pm}$
on $H_L^\pm$ define the sector ground states $\phi_{\chifn}^\pm$
which figure centrally below.

\begin{theorem}[Anti-diagonal sector difference; {\cite{AnalyticKernel}, Thm.~4.1}]
\label{thm:sector-diff}
For an even primitive Dirichlet character $\chifn$ (i.e.,
$\chifn(-1)=+1$, as throughout this paper),
\[
  K_{\chifn}^-(x,y)-K_{\chifn}^+(x,y)
    =-\kappa_{\chifn}(x+y)-K^{\mathrm{arch}}_{\chifn}(x+y).
\]
The primitive part is supported on the anti-diagonals $x+y=\pm m\log
p$ and carries the full character dependence; the archimedean part
depends on $\chifn$ only through the constant $\log(q/\pi)$ (in the
even-character setting; for odd characters the archimedean kernel
acquires an additional $\Re\,\psi(3/4+it/2)$-type term, which is
outside the scope of this atlas).
\end{theorem}

\begin{proof}[Proof sketch]
By Theorem~\ref{thm:kernel},
$K_\chi^\pm(x,y)=\frac{1}{2}[\kappa_\chi(x-y)\pm\kappa_\chi(x+y)]
+\frac{1}{2}[K^{\mathrm{arch}}_\chi(x-y)\pm K^{\mathrm{arch}}_\chi(x+y)]$.
Subtracting $K_\chi^+$ from $K_\chi^-$ cancels all $(x-y)$-arguments
and doubles the $(x+y)$-arguments with sign $-$:
$K_\chi^- - K_\chi^+ = -\kappa_\chi(x+y) - K^{\mathrm{arch}}_\chi(x+y)$.
For the character-dependence claim: $\kappa_\chi$ depends on
$\chifn$ through every $\chifn(p)^m\log p / p^{m/2}$ coefficient
(\eqref{eq:kappa}), whereas
$K^{\mathrm{arch}}_\chi=\mathcal{F}^{-1}[\rho_\chi]$ with
$\rho_\chi(t)=\log(q/\pi)+\Re\,\psi(1/4+it/2)$ depends on $\chifn$
only through the constant $\log(q/\pi)$ (the $\psi$-term is
character-independent in the even-only setting). Full derivation:
\cite{AnalyticKernel}, \S4.1.
\end{proof}

\subsection{The gap functional}

Let $\lambda_1^{\min}(Q_{\chifn}^\pm)$ denote the smallest eigenvalue
of $Q_{\chifn}$ restricted to $H_L^\pm$, and write
\begin{equation}
\label{eq:gap-def}
\begin{aligned}
  \gap_{\chifn}(\lambda)
    &:=\lambda_1^{\min}(Q_{\chifn}^-)
      -\lambda_1^{\min}(Q_{\chifn}^+)\\
    &=\langle\phi_{\chifn}^-|K_{\chifn}^-|\phi_{\chifn}^-\rangle
      -\langle\phi_{\chifn}^+|K_{\chifn}^+|\phi_{\chifn}^+\rangle.
\end{aligned}
\end{equation}
For the trivial character (classical Riemann setting), $\gap_\zeta
(\lambda)\gtrsim\sqrt{\lambda}$ unconditionally \cite{RH_II}. For
non-trivial $\chifn$, the empirical observation is that
$|\gap_{\chifn}|=O(1)$ at the accessible $\lambda$, with
\emph{character-specific} sign; the atlas of
Section~\ref{sec:validation} is the quantitative image of this
observation.

\begin{remark}[Convention-dependence of sign accuracy]
\label{rem:sign-convention-dependence}
The sign accuracy reported in Table~\ref{tab:regression} is computed
relative to the Connes-operator convention $\tilde{A}_\chi$
(Appendix~\ref{app:sign}) used throughout. Under the equivalent
Weil-$Q$ convention $Q_\chi$, all gap signs invert simultaneously
(both predictor and reference), so the relative agreement
predictor-vs.-reference is unchanged. Thus the sign-accuracy
statistic is convention-internal: it tests self-consistency between
predictor and reference within a chosen convention, rather than an
absolute, convention-free sign of $\gap_\chi$. A convention-free
sign would require an external arithmetic identification (e.g.,
from the spectral side via LMFDB), which we do not provide.
\end{remark}

\section{The gap formula: reduced vs.\ unreduced}
\label{sec:formula}

\subsection{The leading-order reduction}
\label{sec:formula-reduced}

Starting from \eqref{eq:gap-def} and
Theorem~\ref{thm:sector-diff}, the gap admits the representation
\begin{equation}
\label{eq:gap-matrix}
\begin{aligned}
  \gap_{\chifn}(\lambda)
    &= \langle\phi_{\chifn}^-|K_{\chifn}^--K_{\chifn}^+|\phi_{\chifn}^-\rangle\\
    &\quad + \langle\phi_{\chifn}^-|K_{\chifn}^+|\phi_{\chifn}^-\rangle
      - \langle\phi_{\chifn}^+|K_{\chifn}^+|\phi_{\chifn}^+\rangle.
\end{aligned}
\end{equation}
If one approximates both ground states by a common function
$\phi_{\chifn}^\pm\approx\phi_\infty$, the second bracket vanishes
and the first simplifies, by Theorem~\ref{thm:sector-diff}, to
\begin{equation}
\label{eq:gap-reduced}
\begin{aligned}
  \gap_{\chifn}(\lambda)
    &\overset{?}{\approx}
      -\iint \phi_\infty(x)\bigl[\kappa_{\chifn}(x+y)
      +K^{\mathrm{arch}}_{\chifn}(x+y)\bigr]\\
    &\quad\times\phi_\infty(y)\,dx\,dy.
\end{aligned}
\end{equation}
Unfolding $\kappa_{\chifn}$ via \eqref{eq:kappa} and using evenness
of $\phi_\infty$ (even sector) delivers the reduced-gap formula
\begin{equation}
\label{eq:gap-reduced-explicit}
\begin{aligned}
  \gap_{\chifn}(\lambda)
    &\overset{?}{\approx}
      2\sum_{\substack{p\nmid q\\m\log p\leq 2L}}
        \frac{\chifn(p)^m\log p}{p^{m/2}}\\
    &\quad\times\Phi_\infty(m\log p)
      \;+\;\text{arch.\ term},
\end{aligned}
\end{equation}
with $\Phi_\infty(t):=(\phi_\infty\ast\phi_\infty)(t)$
(cf.\ \cite{AnalyticKernel}, Cor.~4.4 and \cite{AnalyticGroundstate},
Cor.~6.1). This is the formula one would naturally transfer from the
trivial character case, where the common ground state $\phi_\infty$
is accessible from the Paley-Wiener analysis of \cite{RH_II}.

\subsection{Empirical status of the reduced formula}
\label{sec:formula-reduced-empirical}

Session 6 of the present programme tested \eqref{eq:gap-reduced-explicit}
with sixteen candidate windows $\phi_\infty$: centered Gaussians of
width $\sigma\in\{0.1,0.3,0.5,1.0,2.0\}$, off-center Gaussians, rectangles
of cut-off $T$, $\cos^2$ windows on $[-T,T]$, the raw
$\chifn$-weighted prime sum, and---most decisively---the \emph{empirical}
single-sector ground state
$\phi_\infty:=\phi_{\chifn}^+$ extracted from the $N=200$ Galerkin
diagonalization.

None of the sixteen candidates reached the success criterion
$\geq\!9/10$ sign matches with $R^2\geq 0.8$. The best
leading-order Gaussian window achieves $6/10$ with $R^2=0.15$; the
empirical single-sector ground state achieves $4/10$ with
$R^2=0.02$. Details of the 16-window scan are recorded in
\cite{Woche4Validation}. The crucial observation is that for all
ten characters the empirical single-sector autocorrelation
$\Phi_\infty$ is nearly identical (numerically within $10^{-3}$):
the reduction $\phi_{\chifn}^+\approx\phi_{\chifn}^-\approx
\phi_\infty$ is therefore not merely quantitatively inaccurate, but
erases the character-specific signal.

\begin{remark}[Where the character signal lives]
\label{rem:signal-locus}
The sector-restricted ground states are, to leading order, the same
function: they are both close to the Paley-Wiener ground state of
the diagonal kernel $\kappa_{\chifn}(x-y)+K^{\mathrm{arch}}_{\chifn}
(x-y)$. Character-specific information is stored in the difference
$\phi_{\chifn}^+-\phi_{\chifn}^-$, which is the fine structure
generated by the anti-diagonal term $\kappa_{\chifn}(x+y)+
K^{\mathrm{arch}}_{\chifn}(x+y)$ of Theorem~\ref{thm:sector-diff}.
Any analytic reduction that absorbs $\phi_{\chifn}^+$ and
$\phi_{\chifn}^-$ into a common function loses exactly this signal.
\end{remark}

\subsection{The unreduced formula}
\label{sec:formula-unreduced}

We retain both ground states separately. Define the sector-difference
autocorrelation
\begin{equation}
\label{eq:Delta-def}
  \Deltachi(t)
    := (\phi^+_{\chifn}\!*\phi^+_{\chifn})(t)
      + (\phi^-_{\chifn}\!*\phi^-_{\chifn})(t),
\end{equation}
\emph{Parity-sign note.} The numerical implementation in our scripts
(Appendix~\ref{app:sign}) computes the autocorrelation
$R_{\phi^\pm}(t):=\int\phi^\pm(s)\phi^\pm(s-t)\,ds$ via
\texttt{np.correlate}, then forms
$\Deltachi^{\mathrm{num}}(t):=R_{\phi^+}(t)-R_{\phi^-}(t)$. By the
parity identity $R_{\phi^+}=(\phi^+{*}\phi^+)_{\mathrm{math}}$
(even sector) and $R_{\phi^-}=-(\phi^-{*}\phi^-)_{\mathrm{math}}$
(odd sector), one has
\[
  \Deltachi^{\mathrm{num}}(t)
    =(\phi^+{*}\phi^+)_{\mathrm{math}}(t)+(\phi^-{*}\phi^-)_{\mathrm{math}}(t),
\]
which is precisely the right-hand side of~\eqref{eq:Delta-def}. Thus
$\Deltachi$ as defined in~\eqref{eq:Delta-def} agrees identically with
the script-computed quantity $\Deltachi^{\mathrm{num}}$, and the
coefficient~$-2$ in~\eqref{eq:S-def} is the unique value that makes
the gap identity (Theorem~\ref{thm:closure}) hold. See
Appendix~\ref{app:sign} for the full derivation tracking all six
relevant sign sources.

and the archimedean sector difference
\begin{equation}
\label{eq:arch-diff-def}
  \archdiff_{\chifn}
    := \langle\phi^-_{\chifn},K^{\mathrm{arch}}_-\phi^-_{\chifn}\rangle
     - \langle\phi^+_{\chifn},K^{\mathrm{arch}}_+\phi^+_{\chifn}\rangle,
\end{equation}
where $K^{\mathrm{arch}}_\pm$ are the sector-restricted archimedean
operators (Theorem~\ref{thm:kernel}). Substituting into
\eqref{eq:gap-matrix} and using Theorem~\ref{thm:sector-diff} yields
the unreduced formula stated below.

\begin{theorem}[Unreduced gap formula]
\label{thm:closure}
Let $\chifn$ be a primitive real character of conductor $q$ with
$\chifn(-1)=+1$, and let $\lambda>1$, $L=\log\lambda$. Let
$\phi^\pm_{\chifn}\in H_L^\pm$ be the sector ground states of
$K_{\chifn}^\pm$, each $L^2$-normalized. Then
\begin{equation}
\label{eq:gap-unreduced}
  \gap_{\chifn}(\lambda)
    = S_{\chifn}(\lambda) + \archdiff_{\chifn},
\end{equation}
with
\begin{equation}
\label{eq:S-def}
  S_{\chifn}(\lambda)
    = -2\sum_{\substack{p\nmid q\\m\log p\leq 2L}}
      \frac{\chifn(p)^m\log p}{p^{m/2}}\,\Deltachi(m\log p).
\end{equation}
\end{theorem}

A derivation is given in \cite{AnalyticKernel}, Remark~4.5 (with
sign convention reconciled against \texttt{build\_W}; see
Appendix~\ref{app:sign}). At the level of the \emph{Galerkin
truncation} at basis size $N$, the equation
\eqref{eq:gap-unreduced} also holds exactly:
\begin{equation}
\label{eq:gap-galerkin}
  \gap_{\mathrm{gal}}^{(N)}(\chifn)
    = S_{\chifn}^{(N)}(\lambda) + \archdiff_{\chifn}^{(N)},
\end{equation}
provided $\phi^\pm_{\chifn}$ are taken as the ground-state
eigenvectors of the Galerkin-truncated operators $K_{\chifn}^{\pm,
N}$ at basis size $N$ and $S^{(N)}$, $\archdiff^{(N)}$ are computed
from these finite-$N$ data.

\subsection{Two numerical realizations of \texorpdfstring{$S_{\chifn}^{(N)}$}{S}}
\label{sec:formula-two-realizations}

A subtle point that drives all of what follows is that
\eqref{eq:S-def} admits \emph{two} distinct numerical
implementations from Galerkin data, which agree to within $10^{-3}$
but not identically.
\begin{itemize}
  \item[(R1)] \emph{Spatial autocorrelation.} Evaluate
    $\phi_{\chifn}^\pm$ on a grid, compute $(\phi^\pm
    *\phi^\pm)(t)$ with a numerical correlation, interpolate to
    the prime abscissae $t=m\log p$, and sum.
  \item[(R2)] \emph{Matrix decomposition.} Decompose the Galerkin
    matrix $W^\pm=W^\pm_{\mathrm{arch}}+W^\pm_{\mathrm{prime}}$ and
    form
    $\mathrm{prime\_diff}_{\mathrm{matrix}}
      :=\langle\phi^-,W^-_{\mathrm{prime}}\phi^-\rangle
      -\langle\phi^+,W^+_{\mathrm{prime}}\phi^+\rangle$.
\end{itemize}
Both quantities estimate $S_{\chifn}^{(N)}$; (R1) is obtained via
interpolated autocorrelation, (R2) via exact matrix inner product.
Their discrepancy is at the level of floating-point and
interpolation error, of order $10^{-3}$. For characters with
$|\gap_{\chifn}|\gtrsim 10^{-2}$ this is negligible. For the weakly
signalled characters $\chifn_{21}$, $\chifn_{29}$, $\chifn_{60}$,
with $|\gap|\lesssim 10^{-2}$, it is of the same order as the
signal. Section~\ref{sec:tautology} is devoted to the consequences.

\section{Numerical validation}
\label{sec:validation}

\subsection{Computational setup}
\label{sec:validation-setup}

For each of the ten real characters $\chifn_D$,
$D\in\{5,8,12,13,17,21,24,29,33,60\}$, we construct the Galerkin
matrices
$W^\pm_{\chifn}$ in the Fourier bases of $H_L^\pm$ at cut-off
$\lambda=20000$ ($L=\log\lambda\approx 9.90$). Matrix assembly uses the
explicit-formula coefficients $\chifn(p^m)\log p/p^{m/2}$ for all
prime powers with $m\log p\leq 2L$ and the archimedean operator
$K^{\mathrm{arch}}_{\chifn}$ computed by Fast Fourier transform of
$\rho_{\chifn}$.

Three Galerkin resolutions are used: $N=200$ locally on a desktop,
$N=400$ locally, and $N=600$ on an 8\,GB Hetzner CCX13 cloud
instance (total wall-clock 1.95 hours across all ten characters). The
smallest eigenvalue of each $W^\pm_{\chifn,N}$ is obtained by a
Lanczos inverse iteration with tolerance $10^{-10}$, producing the
sector ground state $\phi^{\pm,(N)}_{\chifn}$ and the Galerkin gap
$\gap_{\mathrm{gal}}^{(N)}(\chifn)
:=\lambda_1^{\min}(W^-_{\chifn,N})-\lambda_1^{\min}(W^+_{\chifn,N})$.

For every character, the quantities $S_{\chifn}^{(N)}$ (via
spatial-autocorrelation realization (R1)) and
$\archdiff_{\chifn}^{(N)}$ are evaluated from the ground states, and
the sum $S_{\chifn}^{(N)}+\archdiff_{\chifn}^{(N)}$ is compared to
$\gap_{\mathrm{gal}}^{(N)}(\chifn)$ and to the high-$N$ Galerkin
reference gap $\gap_{\mathrm{gal}}^{(N_{\mathrm{src}})}(\chifn)$ with
$N_{\mathrm{src}}\in\{400,600\}$ (cf.\ \S\ref{sec:limitation-external}). Scripts
\texttt{arch\_term\_analysis.py} and
\texttt{server\_n600\_full\_analysis.py} record the raw outputs in
\texttt{ARCH\_TERM\_ANALYSIS.json} and
\texttt{ARCH\_TERM\_N600\_SERVER.json} respectively; post-processing
tables are in \texttt{ARCH\_TERM\_ANALYSIS.md} and
\texttt{ARCH\_TERM\_N600\_ANALYSIS.md}.

\subsection{The central result table}
\label{sec:validation-table}

Table~\ref{tab:validation} summarizes the comparison at $\lambda=
20000$ for all ten characters. For each character we list the
high-$N$ Galerkin reference gap
$\gap_{\mathrm{gal}}^{(N_{\mathrm{src}})}$ (with the resolution
$N_{\mathrm{src}}$ from which it is drawn), the Galerkin gap at
$N=200$ and $N=600$, and the predictor $S+\archdiff$ at both
resolutions. Two observations anchor the entire analysis.

First, at $N=600$ the entries labelled
``$\gap_{\mathrm{gal}}^{(N_{\mathrm{src}})}$'' and
``$\gap_{\mathrm{gal}}^{(600)}$'' coincide identically whenever the
reference was itself computed at $N=600$ (the six characters
$\chifn_8,\chifn_{13},\chifn_{21},\chifn_{29},\chifn_{33},\chifn_{60}$);
this is \emph{not} a verification of the formula, but a tautology
exposed in \S\ref{sec:tautology}.

Second, the predictor $S^{(600)}+\archdiff^{(600)}$ achieves high
linear correlation $R^2=0.95$ against the empirical gaps but drops
to $8/10$ sign accuracy because of two sign flips at $\chifn_{21}$
and $\chifn_{29}$.

\subsection{Regression statistics}
\label{sec:validation-regression}

We record all four canonical regression configurations in
Table~\ref{tab:regression}. The pattern is pivotal for what follows.

\begin{table}[t]
\caption{Sign and linear-regression performance of the predictor
  realizations against the high-$N$ Galerkin reference gap
  $\gap_{\mathrm{gal}}^{(N_{\mathrm{src}})}(\chifn_D)$ at $\lambda=20000$,
  $N\in\{200,600\}$. Sign accuracy is computed relative to the
  Connes-convention sign of the high-$N$ Galerkin reference gap
  (sign-convention note below the table).
  $^*$ N=200 statistics are subject to the Galerkin-truncation
  effects documented in Remark~\ref{rem:sw-compression} and
  \S\ref{sec:discussion-slope}; the N=600 row is the operationally
  preferred result.
  Excluding $\chifn_{21}$ (the single sign-reversal outlier; see
  \S\ref{sec:chi21}) yields 9/9 sign accuracy for $\gap_{\mathrm{gal}}$
  at $N=600$ (and 8/9 for $S+\archdiff$), but exclusion is by
  construction---the excluded character is precisely the failure
  case---and therefore not reported as an independent statistic.}
\label{tab:regression}
\centering
\small
\begin{tabular}{@{}lccc@{}}
\toprule
Configuration & $N$ & sign\_ok & $R^2$ \\
\midrule
$\gap_{\mathrm{gal}}$ (all 10)$^*$ & $200$ & $9/10$ & $0.41$ \\
$S+\archdiff$ (all 10)$^*$ & $200$ & $9/10$ & $0.42$ \\
\midrule
$\gap_{\mathrm{gal}}$ (all 10) & $600$ & $10/10$ & $0.99$ \\
$S+\archdiff$ (all 10) & $600$ & $8/10$ & $0.95$ \\
\bottomrule
\end{tabular}

\smallskip
\footnotesize
\textit{Sign-convention note.} Sign accuracy is relative to the
Connes-operator convention $\tilde{A}_\chi$ used throughout
(Appendix~\ref{app:sign}). The opposite convention (Weil-$Q$,
coefficient $+2$) would invert all signs and yield a complementary
sign-accuracy statistic, but the relative agreement
predictor-vs.-reference is unchanged.
\end{table}

At $N=200$, the predictor $S+\archdiff$ reaches $9/10$ sign accuracy
($R^2=0.42$ on the full sample, $R^2=0.80$ on the subsample
excluding $\chifn_{21}$ post hoc---a by-construction statistic, not
reported as an independent test; cf.\ caption of
Table~\ref{tab:regression}). This apparent numerical closure
motivated the ``Weil-kernel closure paper'' title during early
Session~7 of the development. At $N=600$, however, three things
happen simultaneously:
(i) the Galerkin gap $\gap_{\mathrm{gal}}^{(600)}$ becomes
numerically indistinguishable from the empirical reference for six
of ten characters (tautology, \S\ref{sec:tautology});
(ii) $\chifn_{21}$'s Galerkin gap collapses from $+0.282$ to
$-0.004$, confirming it as an $N$-oscillator
(\S\ref{sec:chi21}); and
(iii) the predictor $S+\archdiff^{(600)}$ gains one point on
$R^2$ but loses two points on sign accuracy, now failing both at
$\chifn_{21}$ (where the $N=600$ Galerkin gap is small) and at
$\chifn_{29}$ (where numerical realization (R1) diverges in sign
from realization (R2)).

\begin{remark}[Siegel-Walfisz compression at $N=600$]
\label{rem:sw-compression}
Comparing the $\gap_{\mathrm{gal}}$-columns at $N=200$ and $N=600$
reveals a systematic compression effect: for the strong-positive
character $\chifn_{12}$, $\gap_{\mathrm{gal}}$ falls from
$+0.0325$ to $+0.0115$ (factor $0.35$); for $\chifn_{17}$,
$\chifn_{24}$, $\chifn_{29}$, $\chifn_{60}$ the same
shrinkage factor holds to within $\pm 0.1$. The strong-negative
characters $\chifn_8$, $\chifn_{13}$, $\chifn_{33}$ are
qualitatively stable between $N=200$ and $N=600$ (relative change
$\lesssim 10\%$). This dichotomy is qualitatively reminiscent of the
Siegel-Walfisz upper bound $|\sum_{p\leq x}\chifn(p)\log p|
=o(\sqrt{x})$ applied to the large-prime tail---characters with a
genuine low-conductor signal appear largely determined already at
$N=200$, whereas weakly-signalled characters have their $N=200$
estimates compressed by a factor approximately $0.35$ at $N=600$---but
no theoretical derivation linking the compression factor to a
Siegel-Walfisz constant is offered here. We flag the empirical
observation as an open question; see \S\ref{sec:outlook}, point on
the theoretical derivation of the Siegel-Walfisz compression.
\end{remark}

\begin{table*}[t]
\caption{$N$-convergence of the Galerkin gap and the Weil-kernel predictor
  at $\lambda=20000$ for ten low-conductor real characters. Note that
  $\gap_{\mathrm{gal}}^{(N_{\mathrm{src}})}$ entries with $N_{\mathrm{src}}=600$
  coincide identically with $\gap_{\mathrm{gal}}^{(N=600)}$ because the
  reference values are themselves outputs of the same Galerkin
  diagonalization at $N=600$; this is the tautology discussed in
  \S\ref{sec:tautology}. There is no external, convention-free
  reference here---all gap values originate from the Galerkin method
  itself; cf.\ \S\ref{sec:limitation-external}.
  Sources: \texttt{ARCH\_TERM\_ANALYSIS.md}, \texttt{ARCH\_TERM\_N600\_ANALYSIS.md}.}
\label{tab:validation}
\centering
\small
\begin{tabular}{@{}lrrrrrr@{}}
\toprule
$\chifn$ & $D$ & $\gap_{\mathrm{gal}}^{(N_\mathrm{src})}$ & $\gap_{\mathrm{gal}}^{(200)}$ & $\gap_{\mathrm{gal}}^{(600)}$ & $(S+\archdiff)^{(200)}$ & $(S+\archdiff)^{(600)}$ \\
\midrule
$\chi_5$  &  5 & $+0.01560\ (400)$ & $+0.01429$ & $+0.01603$ & $+3.35\!\cdot\!10^{-2}$ & $+3.37\!\cdot\!10^{-2}$ \\
$\chi_8$  &  8 & $-0.04902\ (600)$ & $-0.04548$ & $-0.04902$ & $-4.22\!\cdot\!10^{-2}$ & $-4.08\!\cdot\!10^{-2}$ \\
$\chi_{12}$ & 12 & $+0.03367\ (400)$ & $+0.03251$ & $+0.01154$ & $+5.07\!\cdot\!10^{-2}$ & $+2.03\!\cdot\!10^{-2}$ \\
$\chi_{13}$ & 13 & $-0.12504\ (600)$ & $-0.12502$ & $-0.12504$ & $-1.15\!\cdot\!10^{-1}$ & $-1.16\!\cdot\!10^{-1}$ \\
$\chi_{17}$ & 17 & $+0.01318\ (400)$ & $+0.00408$ & $+0.00886$ & $+1.67\!\cdot\!10^{-2}$ & $+1.13\!\cdot\!10^{-2}$ \\
$\chi_{21}$ & 21 & $-0.00423\ (600)$ & $\mathbf{+0.28163}$ & $-0.00423$ & $+2.89\!\cdot\!10^{-1}$ & $+1.14\!\cdot\!10^{-2}$\,\textbf{(F)} \\
$\chi_{24}$ & 24 & $+0.01076\ (400)$ & $+0.01632$ & $+0.00866$ & $+1.74\!\cdot\!10^{-2}$ & $+1.86\!\cdot\!10^{-2}$ \\
$\chi_{29}$ & 29 & $+0.01439\ (600)$ & $+0.02372$ & $+0.01439$ & $+2.58\!\cdot\!10^{-2}$ & $-1.16\!\cdot\!10^{-2}$\,\textbf{(F)} \\
$\chi_{33}$ & 33 & $-0.14221\ (600)$ & $-0.14051$ & $-0.14221$ & $-1.31\!\cdot\!10^{-1}$ & $-1.32\!\cdot\!10^{-1}$ \\
$\chi_{60}$ & 60 & $+0.00521\ (600)$ & $+0.12363$ & $+0.00521$ & $+1.27\!\cdot\!10^{-1}$ & $+7.21\!\cdot\!10^{-3}$ \\
\midrule
\multicolumn{4}{l}{\textit{$S+\archdiff$ sign accuracy ($N=200$):}} & & $9/10$ (R²=0.42) & \\
\multicolumn{4}{l}{\textit{$S+\archdiff$ sign accuracy ($N=600$):}} & & & $8/10$ (R²=0.95) \\
\multicolumn{4}{l}{\textit{$\gap_{\mathrm{gal}}$ sign accuracy ($N=600$):}} & & & $10/10$ (R²=0.99, tautology) \\
\bottomrule
\end{tabular}
\end{table*}

\section{Character atlas in four axes}
\label{sec:atlas}

The raw table of Section~\ref{sec:validation} admits a structural
reading along four taxonomic dimensions: two \emph{filter constants}
(parity and root number, both equal to $+1$ across the entire sample)
and two \emph{informative axes} (archimedean weight ratio and gap
signature). We develop each dimension in turn, in order of increasing
informativeness.

\subsection{Parity cartography}
\label{sec:parity}

The parity $\chifn(-1)\in\{+1,-1\}$ decides the archimedean factor of
$\Lambda(s,\chifn)$: even characters use $\Gamma(s/2)$, odd
characters use $\Gamma((s+1)/2)$. In this paper all ten characters
are Kronecker characters $\chifn_D$ with positive discriminant $D$, and
hence all even. This is deliberate: the sector decomposition
$H_L=H_L^+\oplus H_L^-$ used throughout assumes even parity, which
ensures a non-trivial cosine-sine asymmetry of the anti-diagonal
kernel of Theorem~\ref{thm:sector-diff} (odd characters would
require the paired-space formalism
$\chifn\oplus\bar\chifn$ or a tensor with a sign twist; cf.\
\cite{MetaGalerkinBridge}, \S5.5, footnote on Gamma parity). The
even restriction thus fixes the qualitative regime and leaves room
for the three finer classifications below.

\subsection{Root-number classification}
\label{sec:rootnumber}

For every real Dirichlet character the root number $W(\chifn):=
\epsilon_{\chifn}$ is $+1$ by the classical Gauss-sum formula
$\tau(\chifn)=\sqrt{q}$ for primitive real even characters
\cite{IwaniecKowalski}. All ten characters in our atlas thus share the
same root number $W=+1$, and root number is constant across the
table. Consequently root number cannot discriminate gap signs
within our sample. A larger sample extending to complex
characters (\S\ref{sec:outlook}) would exhibit a genuine root-number
axis with $W\in\{+1,-1,\pm i\}$; we flag this as part of the open
outlook rather than as a result of the present atlas.

\subsection{Archimedean-factor layer}
\label{sec:archlayer}

The quantity
$r_{\chifn}:=|\archdiff_{\chifn}^{(200)}|/|S_{\chifn}^{(200)}|$
(ratio of archimedean to prime contribution at $N=200$,
$\lambda=20000$) varies by nearly two orders of magnitude across the
sample. Table~\ref{tab:archratio} records the values.

\begin{table}[t]
\caption{Archimedean weight ratio
  $r_{\chifn}=|\archdiff|/|S|$ at $N=200$, $\lambda=20000$.}
\label{tab:archratio}
\centering\scriptsize
\begin{tabular}{@{}lcccc@{}}
\toprule
$\chifn_D$ & $|S_{\chifn}|$ & $|\archdiff_{\chifn}|$ & $r_{\chifn}$ & Class \\
\midrule
$\chifn_{29}$ & $2.54\!\cdot\!10^{-2}$ & $4.47\!\cdot\!10^{-4}$ & $1.8\%$ & prime-dominated \\
$\chifn_{24}$ & $1.62\!\cdot\!10^{-2}$ & $1.19\!\cdot\!10^{-3}$ & $7.4\%$ & prime-dominated \\
$\chifn_{33}$ & $1.65\!\cdot\!10^{-1}$ & $3.44\!\cdot\!10^{-2}$ & $21\%$ & prime-dominated \\
$\chifn_5$ & $9.06\!\cdot\!10^{-2}$ & $5.71\!\cdot\!10^{-2}$ & $63\%$ & mixed \\
$\chifn_8$ & $1.16\!\cdot\!10^{-1}$ & $7.43\!\cdot\!10^{-2}$ & $64\%$ & mixed \\
$\chifn_{13}$ & $1.02\!\cdot\!10^{-1}$ & $1.33\!\cdot\!10^{-2}$ & $13\%$ & prime-dominated \\
$\chifn_{17}$ & $2.02\!\cdot\!10^{-2}$ & $3.52\!\cdot\!10^{-3}$ & $17\%$ & prime-dominated \\
$\chifn_{12}$ & $6.51\!\cdot\!10^{-2}$ & $1.44\!\cdot\!10^{-2}$ & $22\%$ & prime-dominated \\
$\chifn_{21}$ & $2.31\!\cdot\!10^{-1}$ & $5.77\!\cdot\!10^{-2}$ & $25\%$ & prime-dominated$^\dagger$ \\
$\chifn_{60}$ & $4.95\!\cdot\!10^{-2}$ & $1.76\!\cdot\!10^{-1}$ & $356\%$ & arch-dominated \\
\bottomrule
\end{tabular}

\smallskip
\footnotesize
$^\dagger$ at $N=200$; the effective weights shift at $N=600$ as
the Galerkin ground states settle (cf.\ \S\ref{sec:chi21}).
\end{table}

Three sub-classes emerge:
\begin{itemize}
  \item \emph{Prime-dominated} ($r_{\chifn}\leq 25\%$):
    $\chifn_{12}$, $\chifn_{13}$, $\chifn_{17}$, $\chifn_{21}$,
    $\chifn_{24}$, $\chifn_{29}$, $\chifn_{33}$. The prime sum
    $S_{\chifn}$ determines both magnitude and sign of the gap;
    $\archdiff_{\chifn}$ shifts the quantitative value by no more
    than a quarter. This is the ``natural'' regime in which
    intuition about the Weil kernel carries.
  \item \emph{Mixed} ($r_{\chifn}\in[60\%,70\%]$):
    $\chifn_5$, $\chifn_8$. Primes and archimedean factor are
    comparable; neither alone is a reliable predictor.
  \item \emph{Arch-dominated} ($r_{\chifn}>100\%$):
    $\chifn_{60}$. Here the archimedean sector difference
    \emph{exceeds} the prime contribution by a factor $\sim 4$, and
    $S_{\chifn}$ and $\archdiff_{\chifn}$ even disagree in sign
    ($S=-4.95\!\cdot\!10^{-2}$, $\archdiff=+1.76\!\cdot\!10^{-1}$).
    The reference gap
    $\gap_{\mathrm{gal}}^{(600)}=+0.005$ is near zero and is carried
    entirely by the archimedean term.
\end{itemize}

The sub-class distinction is structural, not just quantitative: for
prime-dominated characters a future analytic control of
$\Deltachi(m\log p)$ would translate directly into gap-sign
prediction, whereas for arch-dominated characters the archimedean
$\psi$-function asymptotics dominate and an independent analysis is
needed.

\begin{remark}[Descriptive vs.\ predictive status of the trichotomy]
\label{rem:trichotomy-descriptive}
With ten characters, the prime/mixed/arch trichotomy of
Table~\ref{tab:archratio} is descriptive of the sample, not
predictive for arbitrary primitive real characters: the threshold
values ($r\leq 25\%$, $r\in[60\%,70\%]$, $r>100\%$) are chosen to
separate the observed clusters, with no datapoint in the gap
intervals $r\in(25\%,60\%)$ and $r\in(70\%,100\%)$. Whether the
trichotomy generalises is a hypothesis to be tested on the
conductor extension $D\leq 100$ (\S\ref{sec:outlook}, point 4); the
present atlas does not assert it as a universal classification.
\end{remark}

\subsection{Gap-signature atlas}
\label{sec:gapsig}

The four classes below organize the ten characters by the
quantitative nature of their gaps, with the high-$N$ Galerkin
reference $\gap_{\mathrm{gal}}^{(N_{\mathrm{src}})}$ read off
Table~\ref{tab:validation}.

\paragraph{Strong-negative prototype.}
$\chifn_{33}$ is the outstanding case:
$\gap_{\mathrm{gal}}^{(600)}=-0.142$, with near-identical values at
$N\in\{200,400,600\}$. Both prime contribution and empirical gap
are consistently negative, and the gap is the largest in absolute
value over the sample. The character is
$\chifn_{33}(n)=\bigl(\tfrac{33}{n}\bigr)$, with factorization
$33=3\cdot 11$; the corresponding Legendre-symbol values at small
primes are $\chifn_{33}(2,5,7,13,17,19,23,29,31,37)=
(1,-1,-1,-1,-1,-1,+1,-1,+1,-1)$, delivering a clear bias towards
$-1$. We designate $\chifn_{33}$ the \emph{strong-negative
prototype} of the atlas.

\paragraph{Strong-positive prototype.}
$\chifn_{12}$ is the counterpart on the positive side:
$\gap_{\mathrm{gal}}^{(400)}=+0.034$, compressed to
$\gap_{\mathrm{gal}}^{(600)}=+0.012$ at
$N=600$ (Remark~\ref{rem:sw-compression}). The primitive character
mod $12$ has conductor a product of two small primes ($q=12=
4\cdot 3$); earlier asymptotic studies \cite{AsymptoticScan} observed
a slope consistent with zero across three orders of magnitude in
$\lambda$ at $N=200$, which we now reinterpret as an $N$-limited
observation (see also \texttt{THEORIE\_CHI12\_KONSTANZ.md}). Within
the present atlas $\chifn_{12}$ is the strong-positive prototype,
albeit at smaller magnitude than $\chifn_{33}$.

\paragraph{Moderate-negative class.}
$\chifn_8$ and $\chifn_{13}$ both exhibit stable negative gaps
($-0.049$ and $-0.125$ respectively) with consistent signs across
$N\in\{200,400,600\}$. They lie between the strong-negative
prototype and the weak-positive cluster.

\paragraph{Weak-positive cluster.}
$\chifn_{17}$, $\chifn_{24}$, $\chifn_{29}$ form a tight cluster of
small positive gaps ($+0.013$, $+0.011$, $+0.014$). Their absolute
values lie within the Siegel-Walfisz compression regime; their
predictor $S+\archdiff$ retains sign at $N=200$ but can flip at
$N=600$ (case of $\chifn_{29}$, realization (R1); cf.\
\S\ref{sec:validation-regression}). This cluster marks the lower
bound of reliable numerical sign resolution for our current
methods.

\paragraph{Arch-dominated near-zero case.}
$\chifn_{60}$ has $\gap_{\mathrm{gal}}^{(600)}=+0.005$ with
$r_{\chifn_{60}}=356\%$ (above). The predictor $S+\archdiff$ happens
to have the correct sign at $N=600$, but not as a consequence of
the prime contribution, which disagrees in sign with the empirical
gap; instead the archimedean term dominates. We flag $\chifn_{60}$
as a distinct regime warranting a separate analysis.

\paragraph{Composite-conductor oscillator.}
$\chifn_{21}$ ($D=21=3\cdot 7$) has
$\gap_{\mathrm{gal}}^{(600)}=-0.004$ but an $N=200$ Galerkin gap of
$+0.282$ --- a sign flip of factor $66$. It is the only character in our sample exhibiting this
behavior, and we treat it in full detail in \S\ref{sec:chi21}.

\smallskip

The four-axis cartography above is the paper's descriptive payoff:
it groups the ten characters into qualitatively distinct regimes
before, and independently of, any predictive formula. In
particular the prime-dominated/mixed/arch-dominated trichotomy
(Table~\ref{tab:archratio}) and the gap-signature classes
(strong-neg, strong-pos, moderate-neg, weak-pos, arch-dominated,
oscillator) together exhibit six sub-classes in ten characters---a
first sign that the ``single-axis transfer'' of the classical
even-dominance regime \cite{RH_II} is too coarse.

\section{The matrix-decomposition tautology}
\label{sec:tautology}

\begin{proposition}[Tautology]
\label{prop:tautology}
Let $\phi^\pm_{\chifn}$ be the ground-state eigenvectors of the Galerkin-truncated
operators $K^\pm_{\chifn}$ at basis size $N$. Then for any decomposition
$W^\pm = W^\pm_{\mathrm{arch}} + W^\pm_{\mathrm{prime}}$ of the Galerkin matrix,
\[
\begin{aligned}
  \gap_{\mathrm{gal}}^{(N)}
    &= \archdiff + \primediff,\\
  \archdiff
    &:= \langle\phi^-, W^-_{\mathrm{arch}}\phi^-\rangle
      - \langle\phi^+, W^+_{\mathrm{arch}}\phi^+\rangle,\\
  \primediff
    &:= \langle\phi^-, W^-_{\mathrm{prime}}\phi^-\rangle
      - \langle\phi^+, W^+_{\mathrm{prime}}\phi^+\rangle.
\end{aligned}
\]
\end{proposition}

\begin{proof}
Linearity of the inner product on the Galerkin space gives
$\langle\phi^\pm, W^\pm\phi^\pm\rangle = \langle\phi^\pm,
W^\pm_{\mathrm{arch}}\phi^\pm\rangle + \langle\phi^\pm,
W^\pm_{\mathrm{prime}}\phi^\pm\rangle$.
Since $\phi^\pm$ are eigenvectors with eigenvalue
$\lambda_1^{\min}(W^\pm)$, the gap
$\gap_{\mathrm{gal}}^{(N)}$ decomposes as the difference of the two
right-hand sides. $\square$
\end{proof}

Although the propositional content is a trivial linearity identity, its
\emph{methodological consequence} is non-trivial: any apparent agreement
between $\gap_{\mathrm{gal}}^{(N)}$ and $S_\chi + \archdiff_\chi$ at fixed
$N$ is structural rather than predictive, not a prediction. When
$\mathrm{prime\_diff}_{\mathrm{matrix}}$ is replaced by the spatial-domain
autocorrelation formula $S_{\chifn}$, a small numerical discrepancy arises
(of order $10^{-3}$, about $1$--$2\%$ of $S$); the ``9/9 sign accuracy at
$N=200$'' reported in preliminary work of the present author admits two
equivalently consistent interpretations: (a) the discrepancy was
\emph{fortuitously small} at low $N$ for the sample of ten characters and
grows at higher $N$ as the predictor's intrinsic limitations become
resolved; (b) at $N=600$ a \emph{conditioning pathology} sets in for
characters with $|\gap|\lesssim 10^{-2}$, where floating-point and
interpolation errors of order $10^{-3}$ become sign-relevant. The two
readings cannot be distinguished without an external reference (cf.\
\S\ref{sec:limitation-external}). In either case, the $N=200$ agreement
is not a structural property of the formula. At $N=600$ the discrepancy
becomes sign-relevant for $\chifn_{29}$ (see Table~\ref{tab:validation}).

\begin{remark}[Why autocorrelation-numerical $S$ fails]
The autocorrelation $(\phi_{\chifn}^+\ast\phi_{\chifn}^+)(t)$ requires
numerical evaluation on a discrete grid, combined with interpolation to
prime abscissae $t=m\log p$. Floating-point and interpolation errors
at the grid scale $\Delta x = 2L/n_{\mathrm{grid}}$ translate into errors
of order $10^{-3}$ in $S$, which is the same order of magnitude as
$|\gap|$ for weakly-signaled characters like $\chi_{21},\chi_{29}$.
\end{remark}

\begin{corollary}[Where predictive power must come from]
\label{cor:paleywiener}
A genuinely predictive Weil-kernel gap formula requires analytic control
of the ground states $\phi^\pm_{\chifn}$ \emph{without} passing through
Galerkin diagonalization. The route is Paley-Wiener asymptotic analysis
(\`a la Connes \cite{Connes}, adapted to the Dirichlet setting) or a
direct spectral identification of $\Deltachi(t)$ from the local data
of $\chifn$ (parity, root number, conductor, arithmetic factor of the
$L$-function).
\end{corollary}

\section{\texorpdfstring{The $\chi_{21}$ $N$-oscillation}{The chi21 N-oscillation}}
\label{sec:chi21}

The character $\chifn_{21}$, with composite conductor $q=21=3\cdot 7$,
is the only entry of Table~\ref{tab:validation} where the Galerkin
gap at $N=200$ and the empirical reference at $N=600$ disagree in
sign. We quantify the discrepancy and attribute it to a
Galerkin-truncation effect that is resolved at high $N$, not to a
defect of the formula.

Table~\ref{tab:chi21} records the $N$-convergence data at
$\lambda=20000$, from the script \texttt{chi21\_n\_convergence.py}.

\begin{table}[h]
\caption{$N$-convergence of $\chifn_{21}$ at $\lambda=20000$.
  The Galerkin gap is flat between $N=200$ and $N=400$ and then
  jumps by more than an order of magnitude between $N=400$ and
  $N=600$. For comparison the strong-positive character
  $\chifn_{12}$ converges monotonically (last column).}
\label{tab:chi21}
\centering\small
\begin{tabular}{@{}cccc@{}}
\toprule
$N$ & $\gap_{\mathrm{gal}}(\chifn_{21})$ & $S+\archdiff$ & $\gap_{\mathrm{gal}}(\chifn_{12})$ \\
\midrule
$200$ & $+0.2816$ & $+0.2888$ & $+0.0325$ \\
$400$ & $+0.2791$ & $+0.2869$ & $+0.0337$ \\
$600$ & $-0.0042$ & $+0.0114$ & $+0.0115$ \\
\bottomrule
\end{tabular}
\end{table}

Two observations structure the analysis.

\begin{itemize}
  \item[(i)] \emph{Step discontinuity.} Between $N=200$ and $N=400$
    the Galerkin gap is numerically stationary
    ($0.2816\to 0.2791$, change $0.9\%$); it then drops by a factor
    of $\sim 66$ in absolute value and flips sign between $N=400$
    and $N=600$. This is not the signature of a smooth
    $N^{-\alpha}$ convergence; the data are consistent with---but
    do not unambiguously prove---a \emph{suspected quasi-degeneracy
    collapse} (full convergence verified only up to $N=600$) in
    which the Galerkin-truncated sector eigenvalues
    $\lambda_1^{\min}(W^\pm_{\chifn_{21},N})$ for $N\in\{200,400\}$
    correspond to a pseudo-ground state distinct from the
    infinite-$N$ ground state, with the true ground state becoming
    representable only once the basis is large enough. A definitive
    diagnosis would require spectral data
    $\lambda_2-\lambda_1$ and eigenvector overlap
    $|\langle\phi^{(N=200)},\phi^{(N=600)}\rangle|$ at intermediate
    $N\in\{500,700,800\}$; we flag this as the most direct further
    verification (cf.\ \S\ref{sec:outlook}, point on $\chifn_{21}$
    spectral diagnosis).
  \item[(ii)] \emph{Predictor tracks the collapse.} The predictor
    $S+\archdiff$ is dominated by the ground-state autocorrelation
    and therefore \emph{also} collapses between $N=400$ and $N=600$,
    from $+0.287$ to $+0.011$. At $N=600$ the predictor no longer
    agrees in sign with the empirical gap, but its magnitude is now
    consistent with the empirical magnitude to within a factor
    $\sim 3$. This confirms that the formula itself is not
    anomalous: once the Galerkin ground state has converged, the
    predictor converges alongside it, with sign determined by the
    near-zero numerical residuals.
\end{itemize}

For contrast, $\chifn_{12}$ (Table~\ref{tab:chi21}, last column)
shows a monotone $N$-convergence with no sign flip:
$+0.0325\to+0.0337\to+0.0115$. This is the generic behaviour on the
stable characters. The discontinuity for $\chifn_{21}$ is therefore
neither a Galerkin artefact endemic to composite-conductor
characters in general (since $\chifn_{12}$, with $q=12=4\cdot 3$,
is composite), nor a systematic boundary of the method, but a
suspected character-specific near-degeneracy whose definitive
spectral analysis (via $\lambda_2-\lambda_1$ scan and eigenvector
overlaps) lies beyond the scope of the present atlas; see
\S\ref{sec:outlook}.

\begin{remark}[Not a formula defect]
\label{rem:chi21-not-defect}
The sign flip in Table~\ref{tab:chi21} is a property of the
Galerkin operator $W^\pm_{\chifn_{21},N}$ at low $N$, not of the
formula \eqref{eq:gap-unreduced}. Since $S+\archdiff$ is defined
in terms of the Galerkin eigenvectors, any spurious eigenvalue at
low $N$ is inherited by the predictor. The right theoretical
object is $\gap_{\mathrm{gal}}^{(600)}(\chifn_{21})=-0.00423$,
and the right question about $\chifn_{21}$ is whether a
Paley-Wiener analysis of the infinite-$N$ ground state would recover
this value directly, without passing through the $N=200,400$
plateau.
\end{remark}

\section{Discussion}
\label{sec:discussion}

\subsection{Implications for the three-component decomposition}
\label{sec:discussion-three-component}

The three-component decomposition \eqref{eq:three-component} divides
any RH-type question into the Galerkin bridge $\mathrm{GBC}$, the
Connes-style operator-identification component $C^{(2)}$, and the
Hurwitz reality-of-zeros transfer. For the trivial character the
even-dominance theorem of \cite{RH_II} closes $C^{(2)}_{X=\zeta}$
unconditionally via the $\sqrt{\lambda}$ separation of the sectors.
For non-trivial real $\chifn$ of small conductor, our atlas demonstrates
three structural facts about $C^{(2)}_{\chifn}$:

\begin{itemize}
  \item[(I1)] The Weil-kernel architecture (Theorems~\ref{thm:kernel}
    and~\ref{thm:sector-diff}) is rigorous and transferable. The
    anti-diagonal sector difference carries the full character
    dependence, as predicted.
  \item[(I2)] At $\lambda=20000$ (a single $\lambda$-value), the gap
    $\gap_{\chifn}(\lambda)$ is bounded in absolute value by
    $0.15$ across all ten characters, with character-specific sign.
    A full $\lambda$-asymptotic ($\Omega(\sqrt{\lambda})$ vs.\
    $O(1)$ vs.\ intermediate growth) cannot be inferred from a
    single $\lambda$-measurement and is left for future work;
    preliminary results for $\chifn_5$ and $\chifn_{12}$ over three
    orders of magnitude in $\lambda$ at $N=200$ are consistent
    with $O(1)$ growth (\cite{AsymptoticScan}, supplementary), but
    a systematic $\lambda$-scan over all ten characters is open.
    The coarse ``universal transfer'' hypothesis ``even dominance
    $\Rightarrow$ even sector wins for all $\chifn$'' is refuted
    at $\lambda=20000$ (since six of ten characters have negative
    or near-zero gap); the $\lambda$-asymptotic refutation requires
    the dedicated $\lambda$-scan above.
  \item[(I3)] Identifying $C^{(2)}_{\chifn}$ requires analytic
    control of the sector ground states $\phi^\pm_{\chifn}$
    themselves, not merely of their eigenvalues. Galerkin
    diagonalization delivers the eigenvalues but, by the tautology
    of Proposition~\ref{prop:tautology}, produces no predictive
    information that does not enter the computation as input.
\end{itemize}

In the meta-programme language of \cite{MetaGalerkinBridge}: for
non-trivial $\chifn$, $C^{(2)}_{\chifn}$ is a genuinely open
component, and its closure is the central technical obstacle to a
Hilbert-P\'olya-style proof of GRH in the Dirichlet case. The
analytic architecture here is ready to accept such a closure; what
is missing is the $\chifn$-specific asymptotic analysis of
$\phi^\pm_{\chifn}$.

After the Zookeeper paper \cite{Zookeeper2026}, this missing
component can be phrased more sharply in CCM language. The Riemann
cell now has an independent Zookeeper kernel package---rank-one
decomposition, Cauchy interlacing, and three-lemma closure---beside
the v2.1 even-dominance package. The natural Dirichlet back-transfer
is therefore Route~D: construct a $\chifn$-twisted CCM operator with
a conductor rank-one term and a $\chifn$-weighted arithmetic
perturbation, then compare its defect norm to the atlas control
quantities. In this reading the Paley-Wiener failures above are not a
dead end; they identify the first concrete stress test for the
$\chifn$-twisted CCM route.

\subsection{Limitation: absence of an external reference}
\label{sec:limitation-external}

All gap values reported in Tables~\ref{tab:validation},
\ref{tab:regression}, and~\ref{tab:archratio} originate from the
same Galerkin diagonalization, evaluated at $N\in\{200,400,600\}$.
There is therefore \emph{no external, convention-free reference} in
this atlas: the high-$N$ values labelled
$\gap_{\mathrm{gal}}^{(N_{\mathrm{src}})}$ are the highest available
Galerkin estimates, not measurements of the true infinite-$N$ gap.
Statistical agreement between the predictor $S_\chi+\archdiff_\chi$
and $\gap_{\mathrm{gal}}^{(N_{\mathrm{src}})}$ is therefore a
self-consistency check within the Galerkin method, not an external
validation. A genuine external reference would require one of:
\begin{itemize}
\item[(i)] mp-arithmetic Galerkin computation at $N\geq 2000$ to
  diagnose finite-precision artefacts;
\item[(ii)] direct spectral evaluation via the LMFDB zeros of
  $L(s,\chifn_D)$;
\item[(iii)] an independent Galerkin scheme using a different basis
  (e.g., Hermite functions instead of Fourier).
\end{itemize}
None of these is provided here; we flag this as a key open
methodological item before any predictive claim about
$\gap_{\chifn}(\lambda)$ at the present $\lambda$ can be made.

\subsection{\texorpdfstring{The slope-$0.72$ regression and its meaning}{The slope-0.72 regression and its meaning}}
\label{sec:discussion-slope}

A secondary quantitative fact exposed by our data is the slope of
the linear regression of $\gap_{\mathrm{gal}}^{(N_{\mathrm{src}})}$
against $S+\archdiff$ at $N=200$:
$\gap_{\mathrm{gal}}^{(N_{\mathrm{src}})}\approx 0.72\cdot
(S+\archdiff)-0.01$
(stable-$9$ subsample, $R^2=0.80$, observational only; cf.\
\S\ref{sec:limitation-external}). We recorded this in Session~7 of
the development with enthusiasm, but subsequent reflection shows
its content to be modest. The slope measures the multiplicative
mismatch between the $N=200$ predictor $S^{(200)}+\archdiff^{(200)}$
and the empirical high-$N$ gap. Under the Siegel-Walfisz
compression of Remark~\ref{rem:sw-compression}, this mismatch
disappears at higher $N$: at $N=600$ the Galerkin gap \emph{is}
the empirical reference (by identity, not by prediction), so the
slope tends to $1$ by construction. The $0.72$ is therefore an
$N=200$-specific Galerkin-truncation constant, not a theoretical
invariant of the formula.

A genuine theoretical analogue of the slope would be the
subleading $N$-dependence of the Galerkin ground states
$\phi^\pm_{\chifn,N}$, controlled by the Paley-Wiener roadmap of
\cite{RH_II} adapted to Dirichlet $L$-functions (cf.\
\cite{AnalyticGroundstate}, \S3, and the REG-$\lambda$-conditioned
rate of \cite{RH_II}, \S6). We flag this as an open quantitative
question of independent interest.

\begin{remark}[Status of N=200 statistics]
\label{rem:n200-status}
The interpretation of the slope-$0.72$ as an $N=200$-specific
truncation constant has methodological consequences for the rest of
the atlas: all statistics reported at $N=200$ in
Tables~\ref{tab:validation}, \ref{tab:regression},
and~\ref{tab:archratio} are subject to the same finite-truncation
effect. Where the $N=600$ value differs from the $N=200$ value, the
$N=600$ value is operationally preferred. The $N=200$ data is
retained for transparency of the convergence behaviour, not as a
final result. In particular, the archimedean weight ratios
$r_\chi$ in Table~\ref{tab:archratio} are computed at $N=200$
because the matrix-decomposition $W=W_{\mathrm{arch}}+W_{\mathrm{prime}}$
is most cleanly resolved at this resolution; their values would
shift modestly at $N=600$, but the prime/mixed/arch trichotomy is
qualitatively preserved (verification: cf.\ supplementary material,
\texttt{ARCH\_TERM\_N600\_ANALYSIS.md}).
\end{remark}

\subsection{Sign convention: rigorous algebraic derivation}
\label{sec:discussion-sign}

The coefficient $-2$ in \eqref{eq:S-def} is determined algebraically
by the structure of the Weil quadratic form (Eq.~\eqref{eq:weil-qf}):
the prime-side contribution carries the factor
$-2\sum_{p,m}w_{p,m}R_\phi(m\log p)$, and the sector gap inherits this
factor via the linearity of the inner product on $H_L^\pm$ and the
sector eigenvector expansion. The full derivation, tracking all six
relevant sign sources (Weil-formula prime-side minus, operator
convention, gap orientation, parity action on basis,
convolution-vs-autocorrelation parity identity, and the real-character
absorption factor), is given in Appendix~\ref{app:sign} and
\cite{AnalyticKernelV2}, Theorem~5.1.

The discovery sequence in our development scripts---an initial
implementation with coefficient $+2$ producing $R\!=\!-0.70$ on the
sample, corrected to $-2$ producing $R\!=\!+0.70$ on the same
sample---reflects an early implementation pitfall in the
autocorrelation parity convention (the sign relation
$R_{\phi^-}=-(\phi^-{*}\phi^-)_{\mathrm{math}}$ was initially handled
incorrectly), not a post-hoc fit of the coefficient to the data. The
rigorous derivation closes the convention. The derivation also yields
two equivalent global conventions: the Connes-operator convention
used throughout the atlas (coefficient $-2$) and the Weil-$Q$
convention (coefficient $+2$), which differ by an overall sign on
the prime part. The atlas does not depend on which convention is
chosen, as long as it is applied consistently; all numerical signs
are consistent with the Connes-operator convention.

\section{Outlook}
\label{sec:outlook}

The atlas delivers a precise diagnosis: predictive closure of
$C^{(2)}_{\chifn}$ via Galerkin numerics is structurally
impossible, and the way forward is an analytic control of the
ground-state difference that bypasses Galerkin entirely. Six
concrete next steps are indicated.

\begin{enumerate}
  \item \textbf{Route~D: $\chifn$-twisted CCM transfer.}
    The first post-Zookeeper task is to construct the Dirichlet
    analogue of the CCM/Zookeeper operator: a $\chifn$-twisted
    Epstein--Mellin transform, a conductor rank-one term, and a
    $\chifn$-weighted arithmetic perturbation. The atlas quantities
    $G_{\chifn}$, $R_{\chifn}$, and $\eta^\pm_{\chifn}$ should then
    be re-read as diagnostics for the corresponding defect norm,
    rather than as direct Paley-Wiener predictors.
  \item \textbf{Paley-Wiener asymptotic analysis of
    $\phi^\pm_{\chifn}$.}
    For the trivial character the asymptotic ground state
    $\phi_\infty$ is known explicitly from the Paley-Wiener
    machinery of \cite{RH_II} (\S5--\S6) under the
    REG-$\lambda$ regularity assumption. The adaptation to
    Dirichlet $L$-functions amounts to identifying the
    $\chifn$-twisted analogue of the relevant Fourier-multiplier
    bound and the rational-function contour to which the
    residue-theorem is applied. Three of the five necessary steps
    transfer mechanically (cf.\ \cite{AnalyticGroundstate},
    Thm.~3.1); the remaining two require a character-specific
    treatment of $L'/L(s,\chifn)$ in the critical strip. This is
    the single highest-leverage technical task.
  \item \textbf{Branch-local Hilbert-P\'olya framework (``NE-B
    Boundary'').}
    Building on \cite{RH_II} (\S4 NE-A, \S5 NE-B) and on the
    Selberg analogue \cite{Selberg1956}, the next major project in the meta-programme
    is a structural theorem identifying \emph{which branches of
    the zeta-function tree admit a Hilbert-P\'olya interpretation,
    and which do not.} The present atlas contributes the first
    quantitative evidence that low-conductor Dirichlet characters
    populate the second class (no branch-local Hilbert-P\'olya
    structure at finite $\lambda$), sharpening the boundary
    theorem that will organize the tree.
  \item \textbf{Conductor extension $D\leq 100$.}
    Ten characters is a small atlas. Extending to all primitive
    real characters with $D\leq 100$ ($\approx 30$ additional
    characters including a few complex conjugates grouped in
    paired spaces) would enable statistical tests of the
    parity/root-number/archimedean classification against a
    distribution of observed gap signs. The Galerkin cost is
    $\mathcal{O}(D\cdot N^3)$; at $N=600$ this is $\sim 10$ hours
    of CCX13 compute per ten characters, a one-weekend job.
  \item \textbf{Paired-space formalism for complex characters.}
    Complex Dirichlet characters come in conjugate pairs
    $\chifn\oplus\bar\chifn$; the Weil quadratic form on
    $L^2[-L,L]\otimes\mathbb{C}^2$ is hermitian but not
    self-adjoint. The analytic framework of \S\ref{sec:setup}
    generalizes, but the sector kernel
    (Theorem~\ref{thm:sector-diff}) acquires a Gauss-sum phase
    (\cite{AnalyticKernel}, Prop.~5.2). A clean exposition is
    a prerequisite for testing the root-number axis
    (\S\ref{sec:rootnumber}) against non-trivial values $W\in
    \{+1,-1,\pm i\}$.
  \item \textbf{Theoretical derivation of the Siegel-Walfisz
    compression.} The empirical dichotomy of
    Remark~\ref{rem:sw-compression}---between characters with
    stable $\gap_{\mathrm{gal}}^{(N=200)}$ and those that are
    compressed by a factor $\sim 0.35$ at $N=600$---asks for a
    theoretical classification in terms of the character-specific
    Siegel-Walfisz constant. A candidate formulation is that
    $\chifn$ with a genuine low-prime bias (e.g.\ $\chifn_{33}$,
    $\chifn_{13}$) has a prime sum that saturates the
    Siegel-Walfisz upper bound in the relevant window and is
    therefore $N$-stable, whereas $\chifn$ with a balanced prime
    sum (e.g.\ $\chifn_{12}$) has a prime sum below the upper
    bound and is therefore compressible.
\end{enumerate}

\appendix

\section{Computational details}
\label{app:computation}

\subsection{Galerkin basis and matrix assembly}
We use the standard Fourier orthonormal basis of $H_L^\pm$:
$\{e_0^+=(2L)^{-1/2}\}\cup\{e_n^+(x)=L^{-1/2}\cos(n\pi x/L)\}_{n\geq 1}$
for the even sector and
$\{e_n^-(x)=L^{-1/2}\sin(n\pi x/L)\}_{n\geq 1}$ for the odd sector.
The Galerkin-truncated operators $W^{\pm}_{\chifn,N}$ at basis size
$N$ are assembled by the function
\texttt{build\_W}$(\chifn,\lambda,\mathrm{sector},N)$ of the script
\texttt{arch\_term\_analysis.py}, which decomposes
$W^\pm = W^\pm_{\mathrm{arch}} + W^\pm_{\mathrm{prime}}$ where
\begin{itemize}
\item[(i)] $W^\pm_{\mathrm{arch}}$ is obtained from the Fourier image
  of $\rho_{\chifn}(t)=\log(q/\pi)+\Re\,\psi(1/4+it/2)$ on a grid of
  spacing $\Delta t=10^{-2}$, truncated at $|t|\leq T_{\mathrm{arch}}=
  30$ (chosen so that the Gaussian ground-state spectrum is contained
  within numerical accuracy).
\item[(ii)] $W^\pm_{\mathrm{prime}}$ is obtained by summing
  $\chifn(p)^m\log p/p^{m/2}$ over all $(p,m)$ with $p\nmid q$ and
  $m\log p\leq 2L$, evaluating the basis-function overlaps
  analytically (for cosine-cosine) or by Simpson integration (for
  sine-sine off-diagonal terms).
\end{itemize}
The smallest eigenvalue and its eigenvector are extracted by
Lanczos inverse iteration from \texttt{scipy.sparse.linalg.eigsh}
with tolerance $10^{-10}$ and \texttt{sigma}$=-\log\lambda$ as a
spectral shift.

\subsection{Autocorrelation evaluation}
The spatial-domain autocorrelation
$(\phi^\pm\ast\phi^\pm)(t)$ is computed on a uniform grid of $n_{
\mathrm{grid}}=4096$ points on $[-2L,2L]$ via \texttt{numpy.correlate}
after zero-padding, then interpolated at the prime abscissae
$t=m\log p$ by cubic \texttt{scipy.interpolate.interp1d}. The
resulting realization of $S_{\chifn}$, called $S_{\mathrm{exact}}$
in our scripts, differs from the matrix-based realization
$\mathrm{prime\_diff}_{\mathrm{matrix}}$ by floating-point and
interpolation errors of order $10^{-3}$ ($1$--$2\%$ of typical
$|S|$-values); this is the source of the $N=600$ sign flips at
$\chifn_{29}$ and $\chifn_{60}$ documented in
\S\ref{sec:validation}.

\subsection{Scripts}
All data in this paper were generated by
\texttt{arch\_term\_analysis.py} (for $N=200$ local),
\texttt{server\_n600\_full\_analysis.py} (for $N=600$ server run,
1.95\,h on Hetzner CCX13), and
\texttt{chi21\_n\_convergence.py} (for Table~\ref{tab:chi21}).
Post-processing is done by \texttt{ground\_state\_difference\_analysis.py}.
JSON outputs are in \texttt{CORE/zoo-mapping/\_results/}; Markdown
summaries in the same directory.

\section{Sign-convention status}
\label{app:sign}

The coefficient $-2$ in \eqref{eq:S-def} is rigorously derived from
the Weil explicit formula in \cite{AnalyticKernelV2},
Theorem~5.1, by tracking all six relevant sign sources (Weil-formula
prime-side minus, operator convention, gap orientation, parity
action on basis, convolution-vs-autocorrelation parity identity,
and the real-character absorption factor). The derivation yields
\emph{two equivalent conventions}:
\begin{enumerate}
  \item the \textbf{Connes-operator convention} $\tilde A_{\chifn}$,
    in which the prime terms are added with positive sign in the
    matrix assembly; this matches the script
    \texttt{build\_W} (i.e.\ $\kappa^{\mathrm{script}}_{\chifn}(u)
    =+\sum w_{p,m}[\delta(u-m\log p)+\delta(u+m\log p)]$) and
    yields coefficient~$-2$;
  \item the \textbf{Weil-$Q$ convention} $Q_{\chifn}$ used in
    \eqref{eq:kappa}, in which $\kappa_{\chifn}$ carries a leading
    minus sign; here the analogous derivation yields
    coefficient~$+2$ (as in \cite{AnalyticKernel}, Cor.~4.4).
\end{enumerate}
The two operators satisfy
$\tilde A_{\chifn}^{\mathrm{prim}}=-Q_{\chifn}^{\mathrm{prim}}$ (the
prime parts differ by an overall sign); both are valid descriptions
of the same spectral problem, but they give oppositely signed gaps.
The atlas and the script consistently use the Connes-operator
convention throughout, so the coefficient~$-2$ in
\eqref{eq:S-def} is the correct one.

A second subtlety concerns the object being summed. The numerical
implementation of \eqref{eq:S-def} (R1, spatial autocorrelation)
uses \texttt{np.correlate}, which computes the \emph{autocorrelation}
$R_{\phi^\pm}(u):=\int\phi^\pm(s)\phi^\pm(s-u)\,ds$.
For the \emph{even} ground state $\phi^+$ one has
$R_{\phi^+}=(\phi^+{*}\phi^+)_{\mathrm{math}}$, but for the
\emph{odd} ground state $\phi^-$ the parity identity gives
$R_{\phi^-}=-(\phi^-{*}\phi^-)_{\mathrm{math}}$ (proof: substitution
$v=-s$ in the integral $\int\phi^-(s)\phi^-(s-u)\,ds$, using
$\phi^-(-v)=-\phi^-(v)$). Consequently, the numerical expression
$\Deltachi^{\mathrm{num}}:=R_{\phi^+}-R_{\phi^-}$ computed by the
script equals
$(\phi^+{*}\phi^+)_{\mathrm{math}}+(\phi^-{*}\phi^-)_{\mathrm{math}}$
in mathematical convolutions---a \emph{sum}, which is precisely the
right-hand side of \eqref{eq:Delta-def}. Thus $\Deltachi$ as defined
in \eqref{eq:Delta-def} agrees identically with the script-computed
quantity $\Deltachi^{\mathrm{num}}$. The coefficient~$-2$ in
\eqref{eq:S-def} is then the unique value that makes the gap
identity (Theorem~\ref{thm:closure}) hold. See
\cite{AnalyticKernelV2}, \S4 and Theorem~5.1, for the full
derivation. All reported signs and all tables
(Tables~\ref{tab:validation}, \ref{tab:regression},
\ref{tab:archratio}, \ref{tab:chi21}) are internally consistent
and have been verified against a direct Galerkin diagonalization.

% --- Bibliography ---
\begin{thebibliography}{99}

% ----- FST master references -----

\bibitem{MathMaster}
L.~Geiger,
\textit{The Zeta Zoo --- The Mathematical Side of Functional
Stability Theory: Branch-Local Hilbert--P\'olya via the Trinity of
UCU, SGE, and Weil-Form Transversality},
Preprint, Zenodo (2026), DOI:
\href{https://doi.org/10.5281/zenodo.19673226}{10.5281/zenodo.19673226}.
Three-component decomposition of $\mathrm{RH}$ in \S2.4; UCU in \S7;
CRM-V classification remark in \S6.5.

\bibitem{RFEPFramework}
L.~Geiger,
\textit{FST Spectrum Duality: The Renormalized Free Energy Principle
as Physical Instantiation of Functional Stability Theory},
Preprint, Zenodo (2026), DOI:
\href{https://doi.org/10.5281/zenodo.19036190}{10.5281/zenodo.19036190}.

\bibitem{Zookeeper2026}
L.~Geiger,
\textit{The Spectral Zookeeper: The Riemann Hypothesis via Spectral
Microcluster Closure in the Connes--Consani--Moscovici Framework},
Preprint (2026), Zenodo DOI
\href{https://doi.org/10.5281/zenodo.19673126}{10.5281/zenodo.19673126}.

% ----- Internal project references (preprint / supplementary notes) -----

\bibitem{RH_II}
L.~Geiger,
\textit{From Landscape to Proof: The Even-Dominance Route to the
Riemann Hypothesis},
Preprint (2026), Zenodo DOI \href{https://doi.org/10.5281/zenodo.19035640}{10.5281/zenodo.19035640}.
Companion: \path{01_RH_TRILOGY/paper/RH_II_Even_Dominance.tex}.

\bibitem{MetaGalerkinBridge}
L.~Geiger,
\textit{A Meta Galerkin Bridge: Diagnostic Decomposition of RH-Type
Programmes}
(Draft 2026), archived predecessor of the \texttt{CORE/zetazoo}
Boundary paper. Content absorbed into \cite{MathMaster} \S2.4.

\bibitem{AnalyticKernel}
L.~Geiger,
\textit{Analytic Kernel: Sector-Kernel Decomposition of the
Dirichlet Weil Quadratic Form},
Supplementary note (2026-04-17, revised 2026-04-16),
\path{CORE/zoo-mapping/ANALYTIC_PIPELINE.md}.

\bibitem{AnalyticKernelV2}
L.~Geiger,
\textit{Analytic Kernel V2 --- Rigorous Weil-Formula Derivation
of the Sign Coefficient},
Supplementary note (2026-04-17),
\path{CORE/zoo-mapping/ANALYTIC_PIPELINE.md}.
Resolves the sign-convention status by deriving the coefficient
rigorously from the Weil explicit formula and tracking all six sign
sources.

\bibitem{AnalyticGroundstate}
L.~Geiger,
\textit{Analytic Ground State: Paley-Wiener Asymptotics for the
Sector Ground States of Dirichlet Weil Quadratic Forms},
Supplementary note (2026-04-17),
\path{CORE/zoo-mapping/ANALYTIC_PIPELINE.md}.

\bibitem{Woche4Validation}
L.~Geiger,
\textit{Week-4 Validation: 16-Window Falsification of the
Leading-Order Gap Formula},
Supplementary note (2026-04-17),
\path{CORE/zoo-mapping/_results/WOCHE4_VALIDATION_RESULT.md}.

\bibitem{AsymptoticScan}
L.~Geiger,
\textit{Asymptotic Scan of the $\chi_5$ and $\chi_{12}$ Sector Gaps
across Three Orders of Magnitude in $\lambda$},
Supplementary note (2026-04-15),
\path{CORE/zoo-mapping/_results/ASYMPTOTIC_SCAN_2026-04-15.md}.

% ----- External references (verified 2026-05-15) -----

\bibitem{IwaniecKowalski}
H.~Iwaniec and E.~Kowalski,
\textit{Analytic Number Theory},
American Mathematical Society Colloquium Publications \textbf{53},
AMS, Providence, RI (2004), ISBN 0-8218-3633-1, 615~pp.

\bibitem{Weil1952}
A.~Weil,
\textit{Sur les ``formules explicites'' de la th\'eorie des nombres
premiers},
Meddelanden fr\aa n Lunds Universitets Matematiska Seminarium
(vol.\ ded.\ \`a Marcel Riesz), 252--265 (1952).

\bibitem{Selberg1956}
A.~Selberg,
\textit{Harmonic Analysis and Discontinuous Groups in Weakly
Symmetric Riemannian Spaces with Applications to Dirichlet Series},
J. Indian Math. Soc. \textbf{20}, 47--87 (1956).

\bibitem{Connes}
A.~Connes,
\textit{Trace formula in noncommutative geometry and the zeros of
the Riemann zeta function},
Selecta Math. (New Series) \textbf{5}, 29--106 (1999),
DOI \href{https://doi.org/10.1007/s000290050042}{10.1007/s000290050042}.

\bibitem{PaleyWiener}
R.~E.~A.~C.~Paley and N.~Wiener,
\textit{Fourier Transforms in the Complex Domain},
American Mathematical Society Colloquium Publications \textbf{19},
AMS, Providence, RI (1934).

\bibitem{LMFDB}
The LMFDB Collaboration,
\textit{The L-Functions and Modular Forms Database},
\url{https://www.lmfdb.org/}. L-function pages accessed 2026-05-16.

\end{thebibliography}

\end{document}
